
\Needspace{12\baselineskip}
\subsubsection*{Naučnoistraživački rad -- \ClanTri}

\textbf{Naslov rada:}
\emph{Enhancing credit card fraud detection: highly imbalanced data case}.\\

\textbf{Autori:}
Dalia Breskuvienė i Gintautas Dzemyda.\\

\textbf{Izvor:}
\emph{Journal of Big Data}, vol. 11, članak 182,
2024. \cite{breskuviene2024fraud}

\begin{table}[H]
\centering
\caption{Sažeti pregled naučnoistraživačkog rada Breskuvienė i Dzemyde (2024)}
\label{tab:breskuviene-pregled}
\small
\renewcommand{\arraystretch}{1.22}
\begin{tabularx}{\linewidth}{@{}p{0.25\linewidth}Y@{}}
\toprule
\textbf{Stavka} & \textbf{Opis} \\
\midrule

Problem &
Selekcija atributa za detekciju prevara kreditnim karticama kada su
podaci izrazito nebalansirani; uobičajene metode selekcije atributa
zanemaruju manjinsku klasu. \\

Skup podataka &
Tri skupa: sintetički skup transakcija (3\,445\,553 instanci,
0,14\,\% prevara), \emph{Sparkov} skup -- spojene datoteke
\texttt{fraudTrain} i \texttt{fraudTest} (1\,852\,394 instanci,
0,52\,\% prevara) i stvarni evropski skup (284\,807 instanci,
0,17\,\% prevara). \\

Korištene metode &
Nova metoda selekcije atributa FID-SOM zasnovana na samoorganizujućim
mapama (SOM), upoređena s univarijatnom selekcijom (F-test, $\chi^2$,
uzajamna informacija), RFE i važnošću atributa iz XGBoosta.
Klasifikatori: Random Forest, XGBoost i CatBoost. \\

Metrike &
F1-mjera, MCC, G-Mean, AUC-PR i AUC-ROC; svaki rezultat je prosjek
pet ponavljanja s različitim \emph{random seed} vrijednostima. \\

Najbolji rezultat &
Na Sparkov skupu: \textbf{XGBoost + FID-SOM} (7 atributa),
F1 \textbf{0,855}, AUC-PR 0,895, AUC-ROC 0,998, MCC 0,857
(bez selekcije atributa F1 je 0,543). \\

\bottomrule
\end{tabularx}
\end{table}

\textbf{Cilj i istraživački problem.}
Autori polaze od toga da su podaci o prevarama nebalansirani, visoko
dimenzionalni (posebno nakon kodiranja kategoričkih atributa) i podložni
promjeni obrazaca kroz vrijeme (\emph{concept drift}). Cilj rada je
pokazati da pažljivo odabran, manji skup atributa poboljšava klasifikaciju
više nego dodatni atributi.

\textbf{Skup podataka i priprema.}
Za nas je posebno važno što je drugi skup u radu isti Sparkov skup,
pri čemu su autori spojili datoteke za treniranje i testiranje i koristili
11 atributa. Umjesto nasumične podjele, podaci su podijeljeni
\emph{hronološki}: najranijih 80\,\% transakcija čini skup za treniranje, a
kasnijih 20\,\% skup za testiranje. Autori naglašavaju da nasumična
podjela vremenski zavisnih podataka daje previše optimistične rezultate.
Kategorički atributi su kodirani \emph{James-Stein} enkoderom, a numerički
normalizovani na raspon [0, 1]. Namjerno nije korišteno uravnotežavanje
klasa ni podešavanje hiperparametara, kako bi se izolovao čist efekat
selekcije atributa.

\textbf{Poređenje eksperimentalnih rezultata.}
Tabela~\ref{tab:breskuviene-rezultati} prikazuje rezultate iz Tabele 6 rada
za Sparkov skup, i to za model bez selekcije atributa (\emph{baseline})
i za predloženu metodu FID-SOM.

\begin{table}[H]
\centering
\caption{Rezultati na Sparkov skupu iz rada Breskuvienė i Dzemyde}
\label{tab:breskuviene-rezultati}
\small
\setlength{\tabcolsep}{4pt}
\renewcommand{\arraystretch}{1.18}
\begin{tabularx}{\linewidth}{@{}Y c c c c c@{}}
\toprule
\textbf{Model} & \textbf{F1} & \textbf{AUC-ROC} & \textbf{AUC-PR} & \textbf{MCC} & \textbf{G-Mean} \\
\midrule
CatBoost -- svi atributi & 0,750 & 0,991 & 0,769 & 0,765 & 0,791 \\
CatBoost -- FID-SOM & 0,828 & 0,997 & 0,866 & 0,831 & 0,870 \\
Random Forest -- svi atributi & 0,622 & 0,967 & 0,767 & 0,660 & 0,682 \\
Random Forest -- FID-SOM & 0,835 & 0,978 & 0,868 & 0,839 & 0,866 \\
XGBoost -- svi atributi & 0,543 & 0,973 & 0,641 & 0,593 & 0,622 \\
\textbf{XGBoost -- FID-SOM} & \textbf{0,855} & \textbf{0,998} & \textbf{0,895} & \textbf{0,857} & \textbf{0,889} \\
\bottomrule
\end{tabularx}
\end{table}

\textbf{Tumačenje rezultata.}
Na Sparkov skupu selekcija atributa ima veoma veliki efekat: F1-mjera
XGBoosta raste s 0,543 na 0,855, a Random Foresta s 0,622 na 0,835.
Istovremeno, AUC-ROC je visok (0,97--0,998) i za modele bez selekcije,
što pokazuje da AUC-ROC na izrazito nebalansiranim podacima slabo
razlikuje dobre i lošije modele. Zato autori stavljaju naglasak na
AUC-PR, MCC i F1-mjeru. Na stvarnom evropskom skupu, čiji su atributi
već dobijeni PCA transformacijom, efekat selekcije je mnogo manji
(F1 0,796 naspram 0,801 za XGBoost). Ukupno je FID-SOM bio najbolja
metoda u 32 od 45 poređenja (71\,\%).

\textbf{Zaključak i značaj za istraživanje.}
Rad pokazuje da su na našem skupu ansambli stabala (XGBoost, CatBoost,
Random Forest) uspješni, ali da njihove performanse snažno zavise od
izbora atributa. To za naš projekat znači da treba ukloniti identifikatore
i nebitne atribute, a kao metrike koristiti F1, AUC-PR i MCC, a ne
samo tačnost i AUC-ROC. Rad upozorava i da nasumična podjela podataka
može precijeniti performanse, pa je hronološka podjela vrijedna razmatranja.
